\chapter{Literature Review}
\label{ch:litreview}

\section{Shortest paths}

Single-source shortest paths in directed graphs with non-negative weights are
classically solved by Dijkstra's algorithm~\cite{Dijkstra1959}, which we use
throughout in its binary-heap form, running in $O((|V|+|E|)\log|V|)$ time.

\section{$k$-shortest simple paths}

Yen's algorithm~\cite{Yen1971} enumerates the $k$ shortest simple $s\to t$
paths in non-decreasing order of weight using $k$ Dijkstra-style searches per
``spur'' along the previously-found paths. Its runtime is
$O(kn(m+n\log n))$, which is polynomial in $k$ but Yen's algorithm is used
here in a slightly non-standard way: we invoke it as an NSP solver by
enumerating paths until one is found whose weight strictly exceeds
$d_G(s,t)$. If the graph has $r$ distinct simple shortest paths from $s$ to
$t$, this requires examining at least $r$ paths before finding one with
strictly greater cost, and $r$ can be exponential in $|V|$. In favorable
inputs Yen is extremely fast, but on adversarial graphs with many alternate
shortest paths it degenerates.

Eppstein~\cite{Eppstein1998} gives an $O((m + n\log n) + k)$ algorithm for
$k$-shortest paths that need not be simple. NSP requires simplicity and so
Eppstein's algorithm does not directly apply.

\section{Next-to-shortest path: positive results}

The NSP problem was introduced in 1996, and Lalgudi and
Papaefthymiou~\cite{LalgudiPapaefthymiou1997} soon after showed it is
NP-complete in directed graphs with non-negative weights (in
particular, when zero weights are permitted). The undirected case admits a
polynomial-time algorithm: Krasikov and Noble~\cite{KrasikovNoble2004}
gave the first, running in $O(n^3 m)$ time, and a sequence of improvements
brought this down to $O(n^3)$, then $O(n^2)$, and finally to
$O(n\log n + m)$ by Wu~\cite{Wu2010}. For \emph{planar} directed graphs
with positive weights, Wu and Wang~\cite{WuWang2015} gave a
polynomial-time algorithm; they also established structural properties of
solutions on general positively weighted digraphs, but no further progress
followed until Chen, Wein and Zhang.

For \emph{general} directed graphs with strictly positive weights, the
problem remained open for nearly three decades. Chen, Wein, and Zhang
\cite{ChenWeinZhang2025} resolved it in November~2025 with an
$O\bigl(|V|^{4}|E|^{3}\log|V|\bigr)$-time algorithm. The algorithm proceeds
by two reductions -- general digraph to $(s,t)$-straight to $(s,t)$-layered
-- followed by a structural enumeration over $6$-tuples of vertices and
edges that, together with a $2$-vertex-disjoint-paths subroutine in a DAG,
suffice to find any NSP if one exists.

\section{Related primitives}

The CWZ algorithm depends critically on a subroutine for finding two
vertex-disjoint paths between prescribed source-sink pairs in a directed
acyclic graph. Fortune, Hopcroft, and Wyllie~\cite{FortuneHopcroftWyllie1980}
gave the first polynomial-time algorithm for the $k$-vertex-disjoint-paths
problem in DAGs for fixed $k$. Tholey~\cite{Tholey2012} subsequently gave a
linear-time algorithm for the special case $k=2$. The CWZ paper cites Tholey
for the $O(|V|+|E|)$ bound on the inner subroutine. In the present work we
implement the simpler Fortune--Hopcroft--Wyllie algorithm for two paths,
which runs in $O(|V|^{2}+|V||E|)$ time (Section~\ref{sec:twovdp}), so our
implementation does not attain the paper's bound.
